\section{Method}
\label{sec:method}

We consider a dataset of cells $\{x_i\}_{i=1}^N$ with $B \geq 1$ batches, where $x_i \in \mathbb{R}^G$ denotes gene expression and $b_i \in \{1, \dots, B\}$ is the batch label. The model consists of a batch-conditional encoder mapping cells to a shared latent space, a set of learnable prototypes representing metacells, and a batch-conditional decoder that reconstructs each cell directly from its own encoding, so that batch-specific variation is handled at decoding rather than absorbed into shared prototypes.

\subsection{Model Components}

\paragraph{Input.}
We assume gene expression profiles $x_i \in \mathbb{R}^G$ with batch labels $b_i$, and a precomputed, sparse cell--cell similarity matrix (graph) $P \in \mathbb{R}^{N \times N}$, approximately block-diagonal by batch (construction and reasoning in Appendix~\ref{app:graph_structure}). By default $P$ is an adaptive Gaussian (RBF) kernel on PCA embeddings~\citep{seacells}, correcting for local density so rare cell types stay well-connected; disentangling the biological and batch-driven variation this mixes together is delegated to the batch-conditional encoder during training. More broadly, $P$ can encode any biologically meaningful similarity---including spatial context---while the reconstruction objective keeps the model grounded in gene expression (Appendix~\ref{app:spatial_affinity}).

\paragraph{Encoder.}
A conditional encoder maps each cell to latent space:
\begin{equation}
z_i = f_\theta(x_i, b_i) \in \mathbb{R}^d.
\end{equation}

\paragraph{Prototypes.}
We introduce $K$ learnable prototypes $\{c_k\}_{k=1}^K$, $c_k \in \mathbb{R}^d$, representing metacells in the latent space.

\paragraph{Soft Assignment.}
Each cell is represented as a soft mixture of prototypes:
\begin{equation}
s_{ik} =
\frac{\exp\left( \frac{z_i^\top c_k}{\tau} \right)}
{\sum_{k'} \exp\left( \frac{z_i^\top c_{k'}}{\tau} \right)},
\end{equation}
where $\tau > 0$ is a temperature parameter controlling assignment sharpness, and both $z_i$ and $c_k$ are $\ell_2$-normalized (see Appendix~\ref{app:training}). Rather than treating $\tau$ as a fixed hyperparameter, we calibrate it from data after pretraining and prototype initialization (see Appendix~\ref{app:calibration}).

The assignment vector $s_i \in \mathbb{R}^K$ encodes the soft membership of cell $i$ across all $K$ prototypes. This representation is used to preserve similarity structure by mapping highly connected cells to similar assignment vectors.

\paragraph{Community Structure Preservation.}
The idea is simple: cells the affinity graph calls related should end up with similar prototype assignments, and cells it calls unrelated should end up with dissimilar ones. We measure assignment similarity on a stop-gradient copy of the prototypes,
\begin{equation}
\hat{s}_{ik} =
\frac{\exp\left( z_i^\top \mathrm{sg}(c_k) / \tau \right)}
{\sum_{k'} \exp\left( z_i^\top \mathrm{sg}(c_{k'}) / \tau \right)},
\qquad
q_{ij} = \hat{s}_i^\top \hat{s}_j \in [0,1],
\end{equation}
where $\mathrm{sg}(\cdot)$ denotes a stop-gradient, so this loss can only move the encoder and never the prototypes directly (letting it move $\{c_k\}$ instead collapses most of them --- Appendix~\ref{app:training}). We match $q_{ij}$ to the graph with binary cross-entropy:
\begin{equation}
\mathcal{L}_{\text{community}}
=
-\sum_{i,j}
\Big[
P_{ij}\log \tilde q_{ij}
+
(1-P_{ij})\log(1-\tilde q_{ij})
\Big],
\end{equation}
with $\tilde q_{ij}$ clipped away from 0 and 1. We sample positive pairs from $P_{ij}$'s real edges, and negative pairs uniformly within the same batch only. Cross-batch alignment is handled instead by $\mathcal{L}_{\text{cVAE}}$ (Total Objective below).

\paragraph{Cohesion.}
We want similar cells to end up sharing a prototype. Concretely, we update the encoder and the prototypes jointly so that a cell's mean similarity to the other members of its own prototype is as high as possible. Since neither $\mathcal{L}_{\text{community}}$ nor $\mathcal{L}_{\text{cVAE}}$ touches the prototypes (previous paragraph; Total Objective below), this is the only term that shapes $\{c_k\}$ directly. We take the objective itself from the classical clustering literature rather than inventing a new one: it is the standard notion of cluster cohesion~\citep{tan2005data}, equivalent to weighted kernel $k$-means / normalized association~\citep{dhillon2004kernel} once membership is soft. It uses the live assignment $s_i$ (not the frozen $\hat{s}_i$ above), so its gradient reaches both the encoder and the prototypes.

Let $\mathrm{sim}(i,j)$ denote a similarity between cells $i$ and $j$, with $\mathrm{sim}(i,i) := 1$ by convention. We use two such similarities in practice: the affinity graph directly, and a Gram-style similarity built from affinity rows that captures shared neighborhoods rather than direct connections (Appendix~\ref{app:cohesion} defines both precisely). For prototype $k$, define its size and the average similarity a cell $i$ has to $k$'s current members:
\begin{equation}
\mathrm{size}_k = \sum_i s_{ik}, \qquad
\overline{\mathrm{sim}}(i,k) = \frac{\sum_j s_{jk}\, \mathrm{sim}(i,j)}{\mathrm{size}_k}.
\end{equation}
Cohesion is the assignment-weighted average of this quantity over all cells:
\begin{equation}
\mathrm{cohesion} = \frac{1}{N}\sum_i \sum_k s_{ik}\, \overline{\mathrm{sim}}(i,k),
\qquad
\mathcal{L}_{\text{cohesion}} = 1 - \mathrm{cohesion},
\end{equation}
so minimizing $\mathcal{L}_{\text{cohesion}}$ is exactly maximizing that mean similarity. Fixing $\mathrm{sim}(i,i):=1$ is what avoids dead prototypes and lets the model use its full capacity. An empty prototype scores exactly 0; without a self term, a prototype that has only just recruited its first member would score 0 too --- indistinguishable from empty, and with no reason for the model to ever grow it further. With $\mathrm{sim}(i,i):=1$, that same prototype scores 1 instead, giving it a foothold to keep recruiting rather than being abandoned outright.

\paragraph{Estimating cohesion from a mini-batch.}
Scoring even one mini-batch cell against $\overline{\mathrm{sim}}(i,k)$ requires $\mathrm{size}_k$ and every other cell's current assignment --- the full column of $s$, not just the mini-batch's. We keep this in a memory: an $N \times K$ table $S^{\text{mem}}$ holding every cell's most recently computed assignment. At each step, the mini-batch's own rows are live (gradient-carrying); every other row is read from the memory as-is. Because a cell can go several steps between visits, we can optionally discount its remembered row by age --- how long it has been since it last appeared in a mini-batch --- so a long-unseen cell counts for less than one just computed (Appendix~\ref{app:cohesion} gives the exact discount and why we leave it off in practice); a cell that reappears simply overwrites its old row with a fresh one. For multi-batch data, cohesion for a cell in batch $b$ is restricted to $V(b)$, the cells either in $b$ or directly connected to it, rather than all $N$ cells (Appendix~\ref{app:cohesion}).

\paragraph{Total Objective.}
\begin{equation}
\mathcal{L} =
\lambda_{\text{comm}}\, \mathcal{L}_{\text{community}}
+
\lambda_{\text{cvae}}\, \mathcal{L}_{\text{cVAE}}
+
\lambda_{\text{coh}}\, \mathcal{L}_{\text{cohesion}},
\end{equation}
where $\mathcal{L}_{\text{cVAE}}$ is the same per-cell reconstruction loss used to pretrain the encoder and decoder (Stage~1 below); we keep optimizing it during joint training, with the batch-embedding parameters learned in Stage~1 frozen. At this point we do not want the model to learn any new batch structure --- the goal is for $\mathcal{L}_{\text{community}}$ and $\mathcal{L}_{\text{cohesion}}$ to shape the encoder and prototypes around correct \emph{within-batch} similarity --- but the decoder still needs $z_i$ to carry enough information to reconstruct real gene expression, so $\mathcal{L}_{\text{cVAE}}$ stays in the objective specifically to preserve the cross-batch alignment Stage~1 already achieved, rather than to teach the encoder anything new about batches. In practice we include two instances of $\mathcal{L}_{\text{cohesion}}$, one for each choice of $\mathrm{sim}$ in Appendix~\ref{app:cohesion}; every term above has its own weight; see Appendix~\ref{app:training} for training details.

\paragraph{Training Procedure.}

\textbf{Stage 1: Pretraining.}
We pretrain the encoder and decoder using scPoli~\citep{scpoli}, a cVAE-based architecture for single-cell data that encodes batch identity via trainable batch embeddings rather than fixed one-hot vectors, enabling effective batch effect correction.
The reconstruction part of pretraining is a per-cell MSE loss on log-normalized gene expression:
\begin{equation}
\mathcal{L}_{\text{cVAE}} =
\frac{1}{N}\sum_i \|x_i - g_\phi(z_i, b_i)\|_2^2,
\end{equation}
where $z_i = f_\theta(x_i, b_i)$. Like any VAE, pretraining also includes a KL term pulling the encoder's approximate posterior toward a standard normal prior, weighted by scPoli's own default schedule. This stage trains the encoder and batch embeddings before prototypes are introduced (see Appendix~\ref{app:training} for architecture and hyperparameter details).

\textbf{Stage 2: Prototype initialization and learning.}
Prototypes are initialized via waypoint initialization~\citep{seacells} from the pretrained encoder embeddings, with temperature $\tau$ calibrated before joint optimization (Appendix~\ref{app:calibration}); batch-embedding parameters from Stage~1 are frozen for the remainder of training (initialization and KL-term details in Appendix~\ref{app:training}).

This two-stage procedure stabilizes training and yields prototypes that preserve within-batch structure, align shared cell states across batches, and represent compact, distinct communities.
Full hyperparameter details, loss weights, and guidance on selecting $K$ are provided in Appendix~\ref{app:training}.